% In C. Verrem Actio Secunda III (in-verrem-2-3) --- English translation
% Translated by Claude (Anthropic) from the Latin (Perseus, Clark text).
% Style guide: STYLE.md.
%
% Book III of the Actio Secunda --- de re frumentaria --- treats Verres's
% administration of the grain tithe in Sicily: the corruption of the lex
% Hieronica, the assessments of the larder, the bagman Apronius, the
% destruction of the farmers province-wide.

\ciceroSection{1} All those, gentlemen, who summon another man into court for the commonwealth's sake --- driven by no enmity, hurt by no private wrong, drawn on by no reward --- ought to weigh not only what burden they take up at present, but how great a business they undertake for their whole life. For they speak as a law unto themselves, of innocence, of self-restraint, of all the virtues, when they call another man to account for his life; and the more so if, as I said before, they do it moved by nothing else save the common good.

\ciceroSection{2} For he who has taken upon himself this work --- to correct the morals of others and to reproach their faults --- who would forgive him if in any matter he himself had departed from the religious obligation of duty? Wherefore such a citizen ought to be praised and loved by all the more, because he not only removes a dishonest citizen from the commonwealth, but also professes himself to be such, and pledges that for him there must be a life lived rightly and honourably --- not only by the common will of virtue and duty, but even by a certain more compelling necessity.

\ciceroSection{3} So this, gentlemen, has often been heard from that most distinguished and most eloquent man, Lucius Crassus: that he repented of nothing so much as of having ever called Gaius Carbo into court. For he had a less free will in all things, and reckoned that his life was being watched by more eyes than he wished. And he --- though fortified by these defences of talent and fortune --- was nevertheless held in by this care, which he had taken upon himself with his judgement not yet firm and at the threshold of life. The virtue and integrity of those who come down to this work as young men is the less observed; that of those who come down with their age now firm, the more. For the first, before they have been able to weigh how much freer is the life of those who have accused no one, accuse for the sake of glory and display; we who have already shown both what we can do and how much we can judge, unless we held our desires easily in check, would never of ourselves cut off from ourselves that licence and freedom of living.

\ciceroSection{4} And I have this greater burden than those who have accused others --- if it should be called a burden, what you bear with gladness and pleasure. Yet I have undertaken this beyond the rest, that it is required of men that they should keep themselves above all from those vices in which they have reproached another. You have accused some thief or some grasping man: every suspicion of greed must always be avoided by you. You have brought up some doer of wrong, or some cruel man: care must always be taken not to seem in any matter the harsher or the more inhuman. A corrupter, an adulterer: it must be diligently provided that no trace of lust appear in your life. In short, all the things you have punished in another must by you yourself be vehemently shunned. For not the accuser only, but not even the rebuker is to be borne, who is reproached for the very vice he reproaches in another.

\ciceroSection{5} I in one man am reproaching all the vices that can be in a ruined and unspeakable man. I say there is no token of lust, of crime, of audacity, that you cannot see in this man's one life. So in this defendant I lay down this law for myself, gentlemen: that I must so live that I may be seen to be most unlike, and always to have been most unlike, this man --- not only in all his deeds and words, but also in that contumacy and arrogance of mouth and eyes that you behold. I bear it, gentlemen, I bear it not heavily, that the life which had before been to me of its own pleasant should now also by my own law and condition be necessary.

\ciceroSection{6} And in this man's case you ask me often, Hortensius, by what enmities or by what wrong I was led down to accuse. I leave aside now the reckoning of my duty and my connection with the Sicilians: of the very enmities I answer you. Or do you think that any enmities are greater than contrary opinions of men, and unlikenesses of pursuits and wills? He who reckons faith most sacred in life --- can he not be an enemy to the man who, as quaestor, dared to plunder, abandon, betray, and assail his consul, with secrets entrusted to him, with money handed over, with all things credited to him? He who cherishes modesty and chastity --- can he with calm mind look upon this man's daily adulteries, his school of harlots, his domestic pandering? He who wishes to keep the religious observances of the immortal gods --- can he not be an enemy to one who has plundered all shrines, who has dared to plunder from the very tracks of the sacred chariots? He who thinks all should be on equal right --- shall he not be most hostile to you, when he reflects on the variousness and caprice of your decrees? He who grieves at the wrongs of allies and the misfortunes of provinces --- shall he not be roused against you by the despoiling of Asia, the harassing of Pamphylia, the squalor and tears of Sicily? He who would have the rights and liberty of Roman citizens held sacred everywhere --- ought he not to be more even than an enemy to you, when he recalls your beatings, when your axes, when your crosses fixed for the punishments of Roman citizens?

\ciceroSection{7} Or if in some matter he had decreed something against my interest unjustly, you would justly think me his enemy: when he has done all things against the affairs, the cause, the reasoning, the interest, and the will of all good men, do you ask why I am his enemy --- whom the Roman people holds for an enemy? I, who especially must take up more even than my common share of burden and duty for the will of the Roman people. What of those things which seem lighter --- can they not move any man's mind: that this man's wickedness and audacity has an easier access to your friendship and the friendship of other great and noble men than any of our virtue and integrity? You hate the industry of new men, you despise their thrift, you contemn their modesty; their talent indeed and virtue you wish pressed down and put out: you love Verres!

\ciceroSection{8} So I suppose: if not by his virtue, not by his industry, not by his innocence, not by his modesty, not by his chastity, then by his conversation, by his learning, by his humanity you are charmed. None of these things has he. On the contrary, all things in him are besmirched with the highest disgrace and shamefulness, with singular stupidity and inhumanity. To this man if anyone's house lies open --- does it lie open, or rather gape and demand something? Him your gatekeepers, him your bedchamber-men love. Him your freedmen, him your slaves and maids love. When he comes, he is summoned out of order; he alone is brought in. Others, often most thrifty men, are shut out. From which can be understood that those are dearest to you who have lived in such a way that without your protection they cannot be safe.

\ciceroSection{9} What? Do you think this is to be borne by anyone --- that we should so live in slender means as to wish to acquire absolutely nothing, that we should defend our own dignity and the kindnesses of the Roman people not by riches but by virtue, while this man, all things snatched from every side, by impunity makes mock, abounds, overflows? That your halls are adorned with his silver, your forum and Comitium with his statues and pictures, especially when you yourselves of your own Mars overflow with all these things? That it is Verres who adorns your villas with his spoils? That it is Verres who vies with Lucius Mummius, that he may seem to have stripped more cities of allies than the other did of enemies; to have adorned more villas with the ornaments of shrines than the other did shrines with the spoils of enemies? And shall he be on this account the dearer to you, that the rest may the more willingly serve your desires at their own peril?

\ciceroSection{10} But these things will be said in another place too, and have been said. Now we shall go forward to the rest, if we shall first have begged a few things of you, gentlemen. Through all the previous speech we have had your minds most attentive: this was very welcome to us. But it will be much more welcome if you wish to attend to what remains, on this account, that in all those things which were said before there was a certain pleasure from the very variety and novelty of the matters and charges. Now we set in hand a case of the grain administration, which by the magnitude of the wrong and by the matter itself will surpass the rest of the charges, but will have less pleasantness in the handling and less variety. It is most worthy of your authority and prudence, gentlemen, in the diligence of hearing to assign no less to scruple than to pleasure.

\ciceroSection{11} In learning this case of the grain administration, set this before yourselves, gentlemen: that you are about to learn of the affairs and fortunes of all the Sicilians, of the goods of those Roman citizens who plough in Sicily, of the revenues handed down by our ancestors, of the life and sustenance of the Roman people. If these things shall seem to you great --- nay greatest --- do not look for how variously and how copiously they shall be spoken of. None of you can fail to see, gentlemen, that all the usefulness and convenience of the province of Sicily, which is joined to the advantages of the Roman people, rests above all in the matter of grain. For in the rest of things we are helped by that province; by this we are fed and sustained.

\ciceroSection{12} The case will fall, gentlemen, into three parts in the prosecution. For first we shall speak about the tithe-grain, then about the bought grain, last about the assessed grain. Between Sicily and the rest of the provinces, gentlemen, in the system of the revenues of fields, this is the difference: that on the rest there is either a fixed tribute imposed (which is called \textit{stipendiary}), as on the Spaniards and on most of the Carthaginians, as it were the prize of victory and the penalty of war; or a censorial leasing has been established, as in Asia by the Sempronian law. The cities of Sicily we received into our friendship and faith on this footing: that they should be on the same right as before, and should obey the Roman people on the same condition on which they had before obeyed their own.

\ciceroSection{13} Very few cities of Sicily have been subdued in war by our ancestors. Of these, the land, although it was made the public property of the Roman people, was nevertheless given back to them; this land is wont to be leased by the censors. There are two cities in alliance whose tithes are not wont to be sold, the Mamertine and the Tauromenitan; and besides these, five cities free of tribute and free, without treaty: Centuripa, Halaesa, Segesta, Halicyae, Panhormus. Besides this all the land of the Sicilian cities is subject to tithe, and so it was, before the empire of the Roman people, by the will and institution of the Sicilians themselves.

\ciceroSection{14} See now the wisdom of our ancestors. When they had joined Sicily, so opportune a support of war and peace, to the commonwealth, they wished to protect and keep the Sicilians with such great care that they not only put no new revenue on their fields, but did not change even the law for selling the tithes, nor either the time or the place of selling: that they should sell at a fixed time of year, in Sicily itself, and finally by the Hieronic law. They wished them themselves to take part in their own affairs, and that their minds should not be moved either by a new law or even by a new name of a law.

\ciceroSection{15} So they decreed that the tithes should always be sold by the Hieronic law, that to them the discharge of that tribute might be the more pleasant: if of the king who was dearest to the Sicilians not only the institutions, with the empire changed, but the very name should remain. By this right the Sicilians always used before Verres became praetor. He first dared to tear down and to change the institutions of all, the custom handed down by the elders, the condition of friendship, the right of partnership.

\ciceroSection{16} In which matter the first thing I reproach and accuse is this: why in a thing so old, so customary, did you do anything new? Did you achieve something by talent? Did you outdo by prudence and counsel so many wisest and most distinguished men, who held that province before you? It is yours, it is the work of your talent and diligence: I grant it and concede it. I know that at Rome, when you were praetor, by your edict you transferred the possession of inheritances from children to strangers, from the first heirs to the second, from the laws to your own caprice. I know that you "corrected" the edicts of all your predecessors and gave the possession of inheritances not to those who produced a will, but to those who said one had been made. And those novelties brought forth and devised by you, I know, were of great gain to you. The same man, I remember, removed and altered the censorial laws too in the matter of letting out repairs of public buildings, that the man whose property it was should not contract for it, that guardians and kinsmen of an orphan should not see to it that he should not be overthrown in all his fortunes; that you should fix a short day for the work, by which to shut the rest out of the business, while in the case of your own contractor you observed no fixed day.

\ciceroSection{17} For which reason I am not surprised that you established a new law in the matter of tithes --- a man so prudent, so practised in praetorian edicts and censorial laws --- I am not, I say, surprised that you devised something. But that of your own accord, without the bidding of the people, without the authority of the senate, you have changed the rights of the province of Sicily, this I reproach, this I accuse.

\ciceroSection{18} In the consulship of Lucius Octavius and Gaius Cotta the senate gave permission that the tithes of wine and oil and of the smaller crops --- which the quaestors had previously been wont to sell in Sicily --- should be sold at Rome, and that for these matters the consuls should lay down such law as seemed good to them. When the contract was being made, the publicans demanded that they should add certain things to the law, but yet should not depart from the rest of the censorial laws. Against this spoke the man who happened then to be at Rome, your guest-friend, Verres --- guest-friend, I say, and familiar of yours --- this Sthenius of Thermae. The consuls examined the case. When they had called many leading and most ample men of the city into council, they pronounced from the council's opinion that they would sell by the Hieronic law.

\ciceroSection{19} Indeed? Most prudent men, endowed with the highest authority, to whom the senate had given full power of laying down laws in the leasing of the revenues, and the Roman people had ordered the same, were unwilling, with one Sicilian protesting, to change the name of the Hieronic law even for the increase of the revenues. And you, a man of the slightest counsel, of no authority, against the bidding of the people and of the senate, with all Sicily refusing, with the greatest loss --- nay, the very ruin --- of the revenues, did away with the whole Hieronic law?

\ciceroSection{20} But what a law is it, gentlemen, that he "corrects" --- nay, that he wholly removes! Most acutely and most diligently written, which law subjected the farmer to all guardianships of the tithe-collector, so that neither in the standing crop nor on the threshing-floors nor in the granaries nor in the moving away nor in the exporting of the grain could the farmer cheat the tithe-collector of one grain without the greatest penalty. The law is so diligently written that it is plain the man who wrote it had no other revenues; so acutely as a Sicilian, so severely as a tyrant. Yet by which law it was profitable to the Sicilians to plough; for the rights of the tithe-collector are so diligently established that nevertheless from an unwilling farmer no more than the tithe could be carried away.

\ciceroSection{21} When these things were so established, Verres in so many years, nay in so many generations, was found to be the man who not to alter but to overturn them, and what had long since been arranged and prepared for the safety of the allies and the advantage of the commonwealth, to convert to his own most dishonest gains. He first appointed certain men under the name of tithe-collectors, in truth the ministers and satellites of his own desires, through whom I shall show that the province was so harassed and laid waste through three years, gentlemen, that for many years we shall not be able to recover it by the innocence and wisdom of many.

\ciceroSection{22} Of all those who were called tithe-collectors the chief was that famous Quintus Apronius whom you see; of whose singular wickedness you have heard the complaints of the gravest embassies. Look, gentlemen, on the man's countenance and look; and from that contumacy which he keeps even in his ruined affairs, consider and recall what you must suppose those Sicilian airs of his to have been. This is the Apronius whom Verres in the whole province, when he had hunted out from every side the most worthless men, and when he had brought with him many like himself, judged most like himself in wickedness, luxury, and audacity. So those men, in a very short time, no business, no reckoning, no recommendation joined together, but the foulness and likeness of their pursuits.

\ciceroSection{23} You know Verres's dishonest and impure ways: feign to yourselves, if you can, someone who in all those things might be his match for the unspeakable lusts of every disgrace. That will be the famous Apronius, who, as he himself signifies not only by his life but by his body and face, is some immense gulf or whirlpool of vices and shames of every kind. This man in all his debauches, this man in his plunderings of shrines, this man in his impure feasts he kept as chief. So great a likeness of ways has its joining and concord that Apronius --- who to others seemed inhuman and barbarous --- to him alone seemed agreeable and witty. So that whom all hated and would not see, without him this man could not be. So that, when others did not even use the same banquets as Apronius, he used the same cups too. Finally so that the most foul smell of Apronius's mouth and body --- which, as men say, not even the beasts could bear --- to him alone seemed sweet and pleasant. He was nearest at the tribunal, alone in the bedchamber, master at the banquet --- and most of all then, when, with the praetor's son in the praetexta reclining beside him, he had begun to dance naked at the banquet.

\ciceroSection{24} This man, as I have begun to say, Verres wished to be chief in harassing and plundering the fortunes of the farmers. Know, gentlemen, that under this praetor the most faithful allies and the best citizens were handed over and made over to his audacity, his wickedness, his cruelty by new institutions and edicts, the whole Hieronic law, as I said before, cast off and repudiated.

\ciceroSection{25} Hear, gentlemen, the famous first edict: \textsc{whatever the tithe-collector shall have declared that the farmer ought to give him for the tithe, that much let the farmer be compelled to give the tithe-collector}. How? "However much Apronius has demanded, let it be given." What is this? Is it a praetor's institution against allies, or the edict and command of a mad tyrant against conquered enemies? "I shall give as much as he has demanded?" He shall demand all that I shall have ploughed. "All? Nay even more," he says, "if he wish." What then? What do you think? Either you shall give it, or you shall be condemned for having acted against the edict. By the immortal gods, what is this? It is not credible.

\ciceroSection{26} I so persuade myself, gentlemen, that, although you may think all things fit this man, yet this seems to you false. For I, when all Sicily was saying it, would yet not dare to affirm it, were I not able to read these edicts in the same words from his own tablets, as I shall do. Give me, please, the clerk, let him read out the proclamation from the codex. Read. \textsc{the edict on declaration}. He says I do not read out the whole. For he seems to signify this by his nod. What do I leave out? Is it that part where you nevertheless take precaution for the Sicilians, and look upon the wretched farmers? For you say that, against the tithe-collector, if he shall have carried off more than was due, you will give an action for eightfold. I do not wish anything to be passed over. Read this too which he asks for, in full. Read. \textsc{the edict on the action for eightfold}. That the farmer should pursue the tithe-collector by an action? Wretched and unjust! You bring men out of the field into the forum, from the plough to the benches, from the practice of country affairs to a strange suit and trial?

\ciceroSection{27} When in all other revenues --- of Asia, Macedonia, Spain, Gaul, Africa, Sardinia, the revenues of Italy itself --- when in all these matters, I say, the publican is wont to be a claimant and a taker of pledges, not a snatcher and possessor, you, against the best, most just, most honourable kind of men --- that is, the farmers --- were establishing those rights which were contrary to all the rest? Which is fairer: that the tithe-collector seek, or that the farmer claim back? That the trial be held with the matter still whole, or with it lost? That he should hold who has sought by his hand, or he who has bid by his finger? What of those who plough with single yokes, who do not themselves leave the work --- in which number, before you were praetor, was a great number and a great multitude of Sicilians --- what shall they do when they have given Apronius what he has demanded? Shall they leave their plough-lands, leave their household Lar? Shall they come to Syracuse, that under you as praetor, of course, on equal right, they may pursue Apronius, your darling and your life, by an action of recovery?

\ciceroSection{28} But suppose so. Some brave and resourceful farmer will be found who, when he has given the tithe-collector as much as the man has said is due, claims it back by an action and pursues the eightfold penalty. I look for the force of the edict, the severity of the praetor: I favour the farmer, I desire Apronius to be condemned for eightfold. What at last does the farmer ask? Nothing but the action for eightfold under the edict. What Apronius? He does not refuse. What the praetor? He orders the recoverers to be challenged. "Let us write the panels." What panels? "From my own staff," he says, "you will challenge." "What? That staff of yours --- of what kind of men is it?" Volusius the soothsayer, and Cornelius the doctor, and those dogs whom you see licking my tribunal. For from the assize he never gave any judge or recoverer; he used to say that all who possessed any clod of land were unjust to the tithe-collectors. One had to come, against Apronius, to those who had not yet breathed out the fumes of an Apronian banquet. O famous and memorable trial! O severe edict! O safe refuge of the farmers!

\ciceroSection{29} And that you may understand of what kind those trials for eightfold were, of what kind those recoverers from his staff were reckoned to be, attend to this. Do you suppose any tithe-collector, this licence having been given that he might carry off from the farmer as much as he had demanded, demanded more than was due? Consider with your own minds whether you suppose any did so, especially since this might happen not from greed only but from imprudence. Many of necessity. But I say all carried off more, and much more, than the tithe. Give me one out of the three years of your praetorship who was condemned for eightfold. Condemned? Nay, against whom an action was even asked under your edict. There was no farmer, of course, who could complain that wrong had been done him; no tithe-collector who confessed that he had had a single grain owed him beyond what was due. Nay, on the contrary, Apronius was snatching and carrying away as much as he wished from each man, and in every place the farmers were complaining, despoiled and harassed; nor will any trial be found.

\ciceroSection{30} What is this? So many brave, honourable, well-favoured men, so many Sicilians, so many Roman knights, hurt by the most worthless and most shameful man, were not pursuing the eightfold penalty, undoubtedly incurred? What cause, what reason is there? That one, gentlemen, which you see: that they saw they would depart from the trial mocked and laughed at into the bargain. For what trial would it be, when from Verres's most shameful and disgraceful retinue three of his hangers-on had sat down under the name of recoverers --- not handed to him by his father but recommended by some little harlot?

\ciceroSection{31} The farmer, of course, would plead his case: would say that no grain had been left him by Apronius; that his goods too had been plundered, that he had been pummelled and beaten. The good men would put their heads together; they would talk among themselves about the carouse and about the little women, if any they could catch as they were leaving the praetor; the matter would seem to be in hand. Apronius would have risen --- a new dignity of a publican; not as a tithe-collector full of squalor and dust, but anointed with unguents, languid with wine and watching --- by his very first motion and breath he would have filled all things with the smell of wine, of unguent, of his body. He would have said the things he was wont to say openly: that he had not bought the tithes, but the goods and fortunes of the farmers; that he, Apronius, was not the tithe-collector, but a second Verres, the master and tyrant of those men. When he had said these things, those most excellent men, the recoverers from his staff, would not have deliberated about acquitting Apronius, but would have asked by what means they could condemn the claimant himself to Apronius.

\ciceroSection{32} When you had given this licence to plunder the farmers to the tithe-collectors, that is, to Apronius --- that he should demand as much as he wished, should carry off as much as he had demanded --- did you prepare this defence for your trial: that you had given an edict that you would give recoverers for the eightfold? Were you, by Hercules, out of the whole resources of the Syracusan assize, of most splendid and honourable men, to give the farmer power not only of challenging but even of choosing the recoverers, yet no one could bear this new kind of wrong: that you, when you had handed over all his crops to the publican and let the matter go from your hands, then should claim back your own goods and pursue them by suit and trial.

\ciceroSection{33} But when in word it is a trial in your edict, in fact really a collusion of your retinue --- of the most worthless men --- with the tithe-collectors, your partners, nay even your agents, do you yet dare to make mention of any trial? Especially when this is refuted not only by my speech but even by the matter itself: that in such great misfortunes of the farmers and wrongs of the tithe-collectors no trial under that famous edict of yours is found, not only carried out but not even asked for.

\ciceroSection{34} Yet he will be milder against the farmers than he seems. For the man who proclaimed by his edict that he would give an action for eightfold against the tithe-collectors, the same man had it in his edict that he would give an action for fourfold against the farmer. Who dares to say this man was hostile or an enemy to the farmers? How much milder he is against them than against the publican! He gave it by edict that what the tithe-collector had declared ought to be given him, that the Sicilian magistrate should exact from the farmer. What is left of the trial that can be given against the farmer? "It is not bad," he says, "that this fear should exist, so that, when it has been exacted from the farmer, the remaining fear of trial may keep him from stirring." If you wish to exact it from me by trial, remove the Sicilian magistrate. If you apply this force, what need is there of trial? Who, moreover, will there be that does not prefer to give your tithe-collectors what they have demanded, rather than be condemned to fourfold by your hangers-on?

\ciceroSection{35} But that famous closing clause of the edict --- that, of all controversies between farmer and tithe-collector, if either wishes, he proclaims he will give recoverers. First, what controversy can there be, when the man who ought to claim takes; and when this man takes not what is owed but what is convenient; while the other from whom it has been taken cannot in any way recover his own by trial? Next, in this matter the muddy man wishes to be cunning and a sharper too, in that he so writes, \textsc{if either wishes, I shall give recoverers}. How prettily he thinks he is stealing! He gives the power to either; but whether he wrote, "if either wishes," or "if the tithe-collector wishes," makes no difference. For the farmer will never wish for those recoverers of yours.

\ciceroSection{36} What of those edicts which, on the spur of the moment, prompted by Apronius, he proclaimed? When Quintus Septicius, a most honourable Roman knight, was resisting Apronius and saying he would give no more than the tithe, there springs up a peculiar sudden edict: \textsc{that no one should remove grain from the threshing-floor before he had come to terms with the tithe-collector}. Septicius bore this iniquity too, and was suffering his grain to be spoiled by rain on the floor, when of a sudden that most fertile and most lucrative edict is born: \textsc{that before the Kalends of August they should have all the tithes carried down to the water}.

\ciceroSection{37} By this edict not the Sicilians (for them indeed he had ruined and afflicted enough by previous edicts), but those very Roman knights, who had reckoned that they could keep their own right against Apronius, splendid men and well-favoured under other praetors, were handed over bound to Apronius. For attend to the kind of edicts they were. \textsc{Let him not remove}, he says, \textsc{from the threshing-floor unless he has come to terms}. This force is great enough to compel an unfair agreement; for I prefer to give more rather than not to remove from the floor in time. But that force does not constrain a Septicius and not a few like Septicius, who say so: "Rather than agree, I shall not remove." To these is opposed: "You must have it carried away before the Kalends of August." "Then I shall carry it away." "Unless you have agreed, you shall not stir it." So a fixed day for carrying away compelled them to remove from the floor; the prohibition of removing, unless he had agreed, applied force to the agreement, not the will.

\ciceroSection{38} Now this further point was not only against the Hieronic law and not only against the custom of his predecessors, but even against all the rights of the Sicilians, which they have from the senate and Roman people: that no man should be compelled to give bail outside his own assize. He laid down that the farmer should give bail to the tithe-collector wherever the tithe-collector wished, so that here too Apronius, when he was binding any man over from Leontini all the way to Lilybaeum, should add a gain of malicious prosecution from the wretched farmers. Although that scheme was found out for malicious prosecution by singular cunning: that he had given by edict that the farmers should declare the acres of their sowings. Which thing, while it had great force for most unjust agreements (as I shall show) and was of no advantage to the commonwealth, was yet for malicious prosecutions, into which all whom Apronius wished should fall.

\ciceroSection{39} For according as anyone had spoken against his will, so an action was demanded against him for the declaration of acres --- by the fear of which trial a great number of bushels of grain was carried off from many men, and great sums of money were exacted; not because it was difficult to declare the number of acres truly --- or even more (for what danger could there be in that?) --- but because there was a cause of demanding a trial, that he had not declared by the edict. Now what trial it was under this praetor, if you remember what staff and what retinue there was, you ought to know. What then is it, gentlemen, that I would have understood from this iniquity of new edicts? That a wrong was done to the allies? You see it. That the authority of his predecessors was repudiated? He will not dare to deny it.

\ciceroSection{40} That Apronius could do so much under this man as praetor? He must confess it. But you perhaps, because the law reminds you, will ask in this place whether he received any money from these things. I shall show that he received the greatest, and I shall convince you that he established all those iniquities of which I spoke before for the sake of his own gain --- if I shall first cast down from his defence that bulwark which he thinks he will use against all my assaults. "I sold the tithes for a great sum," he says. What say you? Did you, most audacious and most insane of men, sell the tithes? Did you sell those parts which the senate and Roman people willed you should sell, or whole crops --- nay, the goods and fortunes of the farmers? If publicly the crier had announced by your command that not the tithes of the grain but the half-shares were being sold, and so buyers had come up to buy half-shares, if you had sold half-shares for more than the rest sold tithes, to whom would it seem wonderful? What? If the crier proclaimed tithes, but in fact --- that is, by law, by edict, by condition --- even more than half-shares were sold, will you yet think this distinguished for you: that you sold for more than was lawful, more than the rest sold for what was right? You sold the tithes for more than the rest.

\ciceroSection{41} By what means did you achieve this? By innocence? Look at the temple of Castor; then, if you dare, make mention of innocence. By diligence? Look at the erasures in your codex on Sthenius of Thermae's entry; then dare to call yourself diligent. By talent? You who in the previous hearing would not question the witnesses, and preferred silently to offer them your face --- speak of yourself and your patrons as men of talent however much you will. By what then have you achieved what you say? Great is the praise if you have outdone your predecessors in counsel, have left to those who follow an example and authority. Perhaps no man was for you fit to imitate; but, as the discoverer and chief of the best practices, all, of course, will imitate you.

\ciceroSection{42} What farmer under you as praetor gave a tithe? Who two? Who did not think himself afflicted with the greatest benefit when he discharged with three tithes for one --- besides those few who, on account of partnership in your thefts, gave nothing at all? See the difference between your harshness and the senate's goodness. The senate, when in the times of the commonwealth it is compelled to decree that second tithes be exacted, decrees that for these tithes money be paid to the farmers, that what is taken beyond what is owed should be reckoned bought, not carried off. You, when you exacted and snatched so many tithes not by senatorial decree but by your new and wicked edicts and institutions, do you reckon you have done a great thing if you have sold for more than Lucius Hortensius (the father of this Quintus Hortensius), or than Gnaeus Pompeius, or than Gaius Marcellus, who did not depart from equity, from law, from custom?

\ciceroSection{43} Or were you to have regard for the reckoning of one year or two, while the safety of the province, the advantage of the grain administration, the system of the commonwealth for posterity were to be neglected? When you had received the matter so established that both enough grain was furnished from Sicily for the Roman people, and yet to the farmers it was profitable to plough and till the fields --- what have you brought about, what have you achieved? That to the Roman people I know not what added to the tithes under you as praetor, you saw to it that the plough-lands were to be deserted and abandoned. Lucius Metellus succeeded you. Are you more innocent than Metellus? More desirous of praise and honour? For you were seeking the consulship, while Metellus was neglecting the honour of his father and grandfather. He sold the tithes for far less, not only than you, but even than those who sold before you. I ask, if he himself could not devise how to sell for the largest sum, could he not even follow your fresh footsteps as the praetor next before him --- so as to use your distinguished edicts and institutions, by you the chief invented and devised?

\ciceroSection{44} He, on the contrary, then thought himself least Metellus if in any matter he should have imitated you. He, from the city of Rome --- a thing no man within human memory has done --- when he thought he must set out to his province, sent letters to the cities of Sicily, by which he urges and asks that they plough, that they sow. As a kindness the praetor seeks this somewhat before his coming, and at the same time shows that he will sell by the Hieronic law --- that is, in all the matter of the tithes will do nothing like this man's doings. And he writes these things, drawn on by no greed --- that he should send letters to another's province before the time --- but with foresight: that, if the time of sowing had passed, we should have no grain from the province of Sicily.

\ciceroSection{45} Hear Metellus's letters. Read. \textsc{Letter of L. Metellus}. These letters, gentlemen, of L. Metellus, which you have heard, ploughed up however much there is of this year's grain from Sicily. No man would have stirred a clod in a tithe-paying field of Sicily, had not Metellus sent this letter. What? Did this come into Metellus's mind by divination? Or was he taught by the Sicilians, who had come to Rome in great numbers, and by the businessmen of Sicily? Of whom, what gatherings to the Marcelli, the most ancient patrons of Sicily, what gatherings to Gnaeus Pompeius (then consul-designate), and to the rest of the connections of that province, were wont to take place, who is ignorant? Which thing indeed, gentlemen, has been done about no man ever: that he should be accused in his absence by those over whose goods and children he held the highest command and power. So great was the force of his wrongs that men preferred to suffer anything rather than not to lament and complain about this man's wickedness.

\ciceroSection{46} Although Metellus had sent these letters to all the cities almost as a suppliant, he could in no part attain the old measure of sowing. For very many had fled away --- a thing I shall show --- and, harassed by this man's wrongs, had left not only the plough-lands but even the seats of their fathers. Not, by Hercules, gentlemen, for the sake of swelling the charge shall I speak; but the impression which I myself received with my eyes and mind, that I shall set out before you truly, and as plainly as I can.

\ciceroSection{47} For when four years later I had come into Sicily, the country seemed to me so afflicted as those lands are wont to be in which a bitter and long-lasting war has been waged. Those plains and most green and shining hills which I had seen before, these I now saw so laid waste and deserted that the field itself seemed to long for its tiller and to mourn its master. The Herbitensian field, the Hennensian, the Murgentine, the Assorine, the Imacharensian, the Agyrinensian had been so largely abandoned that we were searching for the multitude not only of yokes but even of masters. The Aetnensian field, indeed, which had used to be most cultivated, and (which is the head of the grain administration) the Leontine plain --- whose look before was such that, if you had seen it sown, you should not fear a high price of corn --- was so unsightly and rough that in the most fertile part of Sicily we were searching for Sicily. For the year before had vehemently shaken the farmers; the next had overthrown them utterly.

\ciceroSection{48} Do you still dare to make mention to me of the tithes? You, in such great wickedness, in such great bitterness, in so many and such great wrongs --- when in the plough-lands and in the rights of those affairs the province of Sicily stands, the farmers utterly overthrown, the fields abandoned, when in a province so wealthy and so well stocked you have left for any man not the substance, not even any hope --- will you reckon you have something popular when you say you sold the tithes for more than the rest? As if either the Roman people had wished this, or the senate had given you this charge: that, when you had snatched all the fortunes of the farmers under the name of tithes, you should deprive the Roman people of that profit and advantage of the grain administration for the future, and then, if you had added some part of your plunder to the sum of the tithes, you should seem to have deserved well of the commonwealth and well of the Roman people. And I speak so as if his iniquity were to be reproached on this ground: that out of the desire of glory --- to outdo some others by the total of the tithe-grain --- he interposed a more bitter law, harder edicts, and repudiated the authority of all his predecessors.

\ciceroSection{49} "You sold the tithes for a great sum." What? If I show that you diverted home no less than you sent to Rome under the name of tithes, what has your speech of the popular in it, when out of a province of the Roman people you took an equal part for yourself and for the Roman people? What? If I show you carried off twice as much grain as you sent to the Roman people, do we yet think your patron will toss his neck on this charge and give himself up to the people and the crowd? These things you have heard before, gentlemen; but perhaps you so heard them as to have rumour and the talk of all as your authority. Hear now of the boundless money snatched away under the name of grain, that you may at the same time recognize too that dishonest saying of his, who used to say that by the one gain of the tithes he would buy off all his perils.

\ciceroSection{50} We have heard this long since, gentlemen. I deny that any of you has not often heard that the tithe-collectors were this man's partners. I think no other false thing has been said about this man by those who have thought ill of him, save this. For those are to be reckoned partners between whom the matter is shared. I say the whole matter, all the fortunes of the farmers, were this man's own. I say Apronius and Venus's slaves --- which under this praetor was a new kind of publican --- and the rest of the tithe-collectors, were the agents of this man's gain and the ministers of his rapines.

\ciceroSection{51} "How do you show this?" In the same way I showed that this man plundered out of that contract for the columns: from this above all, I think, that he had laid down an unjust and new law. For who has ever attempted to change all rights and the custom of all with reproach and without gain? I shall press on and follow further. You sold by an unjust law, that you might sell for more. Why, the tithes already auctioned and sold --- when nothing more could be added to the total of the tithes, but much to your own gain --- did suddenly and on the spur of the moment new edicts spring up? For that bail should be given to the tithe-collector wherever he wished, that out of the threshing-floor (unless he had agreed) the farmer should not remove, that before the Kalends of August he should have the tithes carried away --- all these things, I say, with the tithes already sold, you proclaimed in the third year. Which, were you doing them for the commonwealth's sake, would have been announced in the selling. Because you were doing them for your own sake, what had been left out by imprudence, this --- prompted by gain and the moment --- you took up.

\ciceroSection{52} But of this whom can you persuade: that without your gain, and without the greatest gain, you should have neglected so great an infamy of yours, so great a peril to your life and fortunes, that, when daily you were hearing the groans and complaints of all Sicily, when, as you yourself said, you reckoned you would be a defendant, when the crisis of this trial was not far from your own opinion, you yet suffered the farmers to be harassed and plundered by the most unworthy wrongs? Surely, although you are of singular cruelty and audacity, you would not wish the whole province to be alienated from you, so many most honourable and most well-off men to become your bitterest enemies, unless the desire of money and that present plunder overcame this reckoning and thought of your own safety.

\ciceroSection{53} For since I cannot give you, gentlemen, the sum and number of the wrongs, while to speak singly of each man's hurt is endless, learn the very kinds of wrongs, please. Nympho is a Centuripine, a busy and industrious man, a most experienced and most diligent farmer. He, when he had taken on great plough-lands as tenant (a thing which even well-off men, as he is, are wont to do in Sicily), and was keeping them at great expense and with a great equipment, was so oppressed by this man's iniquity that he not only abandoned the plough-lands, but even fled from Sicily, and came to Rome together with many cast out by this man. He saw to it that the tithe-collector said Nympho had not declared the number of acres under that famous edict, which had to do with no other thing than gains of this kind.

\ciceroSection{54} When Nympho wished to defend himself by a fair trial, this man gives him most excellent recoverers --- the same Cornelius the doctor (that is Artemidorus of Perga, who in his own home country was once Verres's leader and master in plundering Diana's temple); and Volusius the soothsayer; and Valerius the crier. Nympho is condemned before he has even fairly stood. For how much, perhaps you ask. There was no fixed penalty in the edict: of all that grain which was on the floors. So Apronius the tithe-collector takes from Nympho's plough-lands not the tithe owed, not grain stored away and concealed, but seven thousand medimni of wheat as the penalty of the edict, by no right of the contract.

\ciceroSection{55} The estate of the wife of Xeno of Menae, a most noble man, had been let out to a tenant farmer. The tenant, because he could not bear the wrongs of the tithe-collectors, had fled from the field. Verres was giving against Xeno that condemnatory action of his on the declaration of acres. Xeno denied that the matter concerned him; he said the estate had been let out. This man was giving the action: \textsc{if it should appear that the acres of that estate were more than the tenant had declared}, then Xeno should be condemned. Xeno said that not only had he himself not ploughed --- which was enough --- but that he was neither the owner nor the lessor; that the estate was his wife's; that she herself transacted her own business, that she herself had let it out. Marcus Cossutius, a man of the highest splendour and of the highest authority, was defending Xeno. None the less this man kept giving the action --- 50,000 sesterces. Although that man saw recoverers from the staff of brigands being prepared for him, he yet said he would accept the trial. Then this man, with a great voice, commands Venus's slaves, while Xeno was hearing, to keep the man under guard while the matter was being judged; when judged, to bring him before him; and at the same time he says that he does not think him --- if for his riches he despised the penalty of condemnation --- about to despise even the rods. By this force and this fear he was led to give the tithe-collectors as much as this man commanded.

\ciceroSection{56} Polemarchus is a Murgentine, a good and honourable man. When upon him for fifty acres seven hundred medimni were demanded as the tithe, because he refused, he was led to court at this man's house, and, although this man was even still in bed, he was led into the bedchamber --- which was open to no one save a woman and a tithe-collector. There, when he had been beaten with fists and feet, the man who had not wished to settle for seven hundred medimni promised a thousand. Eubulides Grospus is a Centuripine, a man chief at home both for virtue and nobility, and now too for money. Know, gentlemen, that to this man --- a most honourable man of a most honourable city --- there was left not only of the grain but even of life and blood as much as Apronius's lust permitted. For by force, by trouble, by blows he was driven to give of grain not as much as he owed, but as much as he was being compelled to give.

\ciceroSection{57} Sostratus and Numenius and Nymphodorus, three brothers in partnership of the same city, when they had fled from their fields because more grain was being demanded of them than they had ploughed up, Apronius came with men gathered to their plough-lands, plundered all the equipment, led off the household slaves, drove off the cattle. Afterwards, when Nymphodorus had come to him at Aetna and was begging that his own things be restored, he ordered the man to be seized and hung up on a certain wild olive --- which is a tree, gentlemen, in the forum at Aetna. So long did he hang on the tree, an ally and friend of the Roman people in the city and forum of allies, your settler and farmer, as Apronius's pleasure carried.

\ciceroSection{58} Kinds, gentlemen, of countless wrongs I have for some while been bringing forward, naming them one by one. The boundless multitude of wrongs I leave aside. Set before your own eyes and minds these onsets of the tithe-collectors throughout all Sicily, the plunderings of the farmers, the harshness of this man, the kingship of Apronius. He despised the Sicilians; he reckoned that they would not themselves be vehement in pursuing the matter, and thought that you would bear their wrongs lightly.

\ciceroSection{59} Be it so: he had a false opinion about them, an evil opinion about you. Yet, although he deserved ill of the Sicilians, did he cherish Roman citizens? Did he indulge them? Was he given over to their goodwill and favour? He, the Roman citizens? But to none was he more an enemy or more hostile. I leave aside the chains; I leave aside the prison; I leave aside the floggings, I leave aside the axes; finally I pass over that famous cross, which this man wished to be a witness to Roman citizens of his humanity and goodwill toward them. I leave aside, I say, all these things, and reserve them for another time of speaking. About the tithes, about the condition of Roman citizens in their plough-lands I am disputing. How they were treated, gentlemen, you have heard from themselves: they have said their goods were snatched from them.

\ciceroSection{60} But these things, since the case has been such, are to be borne: that equity has had no force, that custom has had none. Finally, gentlemen, no losses are so great that brave men, endowed with a great and free mind, do not reckon them to be borne. What if hands were laid by Apronius, with no hesitation, on Roman knights not obscure or unknown, but honourable and illustrious? What do you await? What more do you think I should say? Or that we should so handle it that we may make an end with this man the more quickly, that we may the more speedily reach Apronius --- which is what I myself promised him in Sicily? He kept Gaius Matrinius, gentlemen, a man of the highest virtue, the highest industry, the highest favour, in public for two days at Leontini. By Quintus Apronius, gentlemen --- a man born in disgrace, brought up to shamefulness, fitted to Verres's scandals and lusts --- a Roman knight, know, was forbidden food and shelter for two days, kept and watched at Leontini in the forum under Apronius's guards for two days, and not let go before he was beaten down to the man's terms.

\ciceroSection{61} For why should I speak, gentlemen, of Quintus Lollius, an esteemed and honourable Roman knight? The thing I shall tell is famous, well-known throughout all Sicily. He, when he was ploughing in the Aetnensian field, and that field had been handed over to Apronius along with the rest, trusting in his old equestrian authority and favour, declared that he would give the tithe-collectors no more than was owed. His talk is reported to Apronius. Indeed this fellow began to laugh and to wonder that Lollius had heard nothing of Matrinius, nothing of the rest. He sends Venus's slaves to the man. Mark this too: that the tithe-collector had attendants assigned by the praetor, if this can seem any small proof that this man under the name of tithe-collectors abused them for his own gain. Lollius is led, nay dragged, by the Venusians, just as Apronius had returned from the wrestling-school and had reclined in the dining-room he had spread in the forum at Aetna. Lollius is set in that timely banquet of gladiators.

\ciceroSection{62} By Hercules, I should not believe these things which I am saying, gentlemen, although I had heard them commonly, were it not that the old man himself spoke to me with the highest authority, when, weeping, he gave thanks to me and to this purpose of my prosecution. There is set, as I say, a Roman knight nearly ninety years old in Apronius's banquet, while meanwhile Apronius rubbed his head and face with unguent. "What is it, Lollius?" he says. "Do you not know how to do right unless compelled by trouble?" The man, what he should do, whether to be silent or to answer, what at last he should do at that age and with that authority, did not know. Apronius meanwhile was demanding dinner and cups; and his slaves --- who were of the same character as their master, and born of the same kind and place --- were carrying all these things before Lollius's eyes. The guests laughed; Apronius himself roared with laughter --- unless perhaps you suppose he did not laugh in wine and play, who now in his own peril and ruin cannot keep down his laughter. To be brief, gentlemen: by these insults know that Quintus Lollius was compelled to come to Apronius's terms and conditions.

\ciceroSection{63} Lollius, hindered by age and disease, could not come to give testimony. What need is there of Lollius? No one is ignorant of this; no one of your friends, no one brought forward by you, no one questioned by you will say now that he hears this for the first time. Marcus Lollius, his son, a most choice young man, is here. You shall hear his words. For Quintus Lollius, his other son, who accused Calidius --- a young man both good and brave and especially eloquent --- when, moved by these wrongs and insults, he had set out for Sicily, was killed on the journey. Of his death the cause is laid on runaway slaves; in fact no one in Sicily doubts that he was killed because he could not keep his counsels about Verres shut up. This man, moreover, was not in doubt that he, who had before this prosecuted another man led on by zeal, would now be at hand to meet him --- moved by his father's wrongs and his domestic grief.

\ciceroSection{64} Do you now understand, gentlemen, what plague, what monstrousness has been let loose in your most ancient, most faithful, most neighbouring province? Do you now see for what cause Sicily, having before suffered so many men's thefts, rapines, iniquities, and ignominies, could not bear this new and singular and incredible kind of wrongs and insults? Now all understand why the whole province sought as defender of its safety the man from whose faith, diligence, perseverance this man could in no way be torn. You have been present at so many trials, you know that so many guilty and dishonest men have been accused both within your own and your predecessors' memory: have you ever seen, have you ever heard of any man's having been concerned in such great thefts, in thefts so open, in such audacity, in such shamelessness?

\ciceroSection{65} Apronius had Venus's slaves about him as bodyguards; he led them round the cities; he ordered banquets to be prepared for himself publicly, dining-couches to be spread, and to be spread in the forum. Thither most honourable men were summoned --- not only Sicilians but even Roman knights --- so that the most respected and honourable men were held at his banquet, the man with whom no one had ever wished to live save the foul and impure. These things, when you knew them, when you heard them daily, when you saw them, most ruined and most lost of all mortals --- if they were happening without your greatest gain, would you have suffered and granted them to happen at such great peril of your own? Did Apronius's gain, did his most filthy talk and his disgraceful flatteries weigh so much with you that no care or thought of your own fortunes ever touched your mind?

\ciceroSection{66} You see, gentlemen, what and how great a fire from the assault of the tithe-collectors has spread under this man as praetor not only through the fields but even through the rest of the fortunes of the farmers, and not only through their goods but even through their rights of liberty and citizenship. You see some hanging from a tree, others being beaten, others being scourged, others again kept under guard in public, others set up as a spectacle at a banquet, others condemned by the doctor and the crier of the praetor; meanwhile, none the less, the goods of all those men carried off and plundered out of the fields. What is this? The empire of the Roman people? Praetorian laws, courts? Faithful allies, a province at our suburb? Are not all things rather such that, if Athenion the king of the runaway slaves had won, he would not have done them in Sicily? Not, I say, gentlemen, would the insolence of the runaways have attained any part of this man's wickedness. So privately. What? Publicly --- in what way were the cities treated? You have heard, gentlemen, very many testimonies of the cities, and you shall hear of the rest.

\ciceroSection{67} And first hear briefly about the people of Agyrium, faithful and illustrious. Agyrium is among the chief honourable cities of Sicily, of men, before this praetor, well-off and the highest farmers. When the same Apronius had bought the tithes of its field, he came to Agyrium. Having come there with attendants and with force and threats, he began to demand a great sum of money, that he might depart with the gain received: he said he wished to have nothing of the business at all, but, the money received, to run on into another city as soon as possible. All the Sicilians are not to be despised, if our magistrates allow it; but they are men brave enough and downright frugal and sober. And among the chief is this city, gentlemen, of which I am speaking.

\ciceroSection{68} So the Agyrines answer that most dishonest man that they would give him the tithes in such manner as they owed: gain --- since he had bought for a great sum --- they would not add. Apronius makes plain to that man whose business it was. At once, just as if some conspiracy at Agyrium against the commonwealth had been made, or a praetor's legate had been beaten, so the magistrates and chief five from Agyrium are summoned at this man's call. They come to Syracuse. Apronius is at hand. He says that those very men who had come had acted against the praetor's edict. They asked, what? He answered that he would speak to the recoverers. This most just man was casting that fear of his into the wretched Agyrines: he was threatening that he would give recoverers from his own staff. The Agyrines, the bravest of men, said they would suffer the trial.

\ciceroSection{69} This man pressed on Artemidorus Cornelius the doctor and Tlepolemus Cornelius the painter and recoverers of that kind, of whom no one was a Roman citizen, but Greeks, sacrilegious and dishonest from of old, but suddenly Cornelii. The Agyrines saw that whatever Apronius brought before those recoverers, he would prove most easily. They preferred to be condemned with this man's hatred and infamy than to come to his terms and bargains. They asked in what words he was giving the recoverers their charge. He answered, \textsc{if it shall appear that the man has acted against the edict}; in which matter he said he would speak in the trial. They preferred to do battle on the most unjust terms, with the most dishonest recoverers, than to settle anything by their own will with this man. This man sent under Timarchides to warn them, if they were wise, to settle. They flatly refused. "What then? Do you prefer to be condemned to fifty thousand sesterces apiece?" They said they would prefer it. Then this man, with everyone hearing clearly, "He who shall be condemned," he says, "shall be flogged with rods to death." Here they began with weeping to ask and beg that they might be allowed to hand over to Apronius their crops, all their fruits, their plough-lands empty; that they themselves might depart without ignominy and trouble.

\ciceroSection{70} On this condition, gentlemen, Verres sold the tithes. Let Hortensius, if he wishes, say that Verres sold the tithes for a great sum. This was the condition of the farmers under this praetor: that they reckoned matters had gone splendidly with them if they were allowed to hand over the empty fields to Apronius; for they were eager to escape the many crosses set out for them. Whatever Apronius had declared was owed, that much by the edict had to be given. Even if he had declared more than had been grown? Yes, since the magistrates by his edict had to exact it. But the farmer could claim it back. With Artemidorus as recoverer. What if the farmer had given less than Apronius had demanded? An action against the farmer for fourfold. From what panel of judges? From the praetor's distinguished staff of most honourable men. What further? "I say you have declared too few acres: challenge the recoverers, since you have acted against the edict." From what number? From the same staff. What will be the end? "If you shall be condemned, and indeed when you shall be condemned" --- for what doubt of condemnation could there be with those recoverers? --- "you must of necessity be flogged with rods to death." Under these laws, on these terms, will any man be so foolish as to think the tithes were sold, who supposes that nine parts were left to the farmer, who does not understand that this man held for himself, as gain and plunder, the goods, the possessions, the fortunes of the farmers? In fear of the rods the Agyrines said they would do whatever was commanded.

\ciceroSection{71} Hear now what he commanded, and pretend, if you can, that you do not understand that the praetor himself --- a thing all Sicily saw through --- was the buyer of the tithes, nay the master and king of the farmers. He commands the Agyrines that they themselves take the tithes publicly, and give Apronius the gain. If he had bought for a great sum --- since you are the one who most diligently inquired into prices, who, as you say, sold for a great sum --- why did you reckon that gain ought to be added for the buyer? Be it so, you reckoned: for what reason did you command that they add? What else is to take and to gather money (in which the law holds you), if not this --- by force and command to compel unwilling men to give gain to another, that is, to give money?

\ciceroSection{72} What then? They were ordered to give Apronius, the praetor's darling, some little gain. Reckon it given to Apronius, if "Apronius's little gain" and not "praetorian plunder" shall seem so to you. You command that they take the tithes, that they give Apronius as gain 33,000 medimni of wheat. What is this? One city out of one field is being compelled, by the praetor's command, to give to Apronius almost the monthly rations of the Roman commons. Did you sell the tithes for a great sum, when so much gain was given to the tithe-collector? Surely, had you inquired diligently into the price when you were selling, they would have added ten thousand medimni rather than 600,000 sesterces afterwards. The plunder seems great: hear the rest, and listen diligently, that you may wonder less that the Sicilians, compelled by necessity, sought help from their patrons, from the consuls, from the senate, from the laws, from the courts.

\ciceroSection{73} That Apronius might approve this wheat which was given him, Verres commands the Agyrines that for each medimnum a sesterce be given to Apronius. What is this? When such a number of grain has been commanded and wrung out under the name of gain, are coins demanded besides for the wheat to be approved? Or could not Apronius --- nay, anyone --- if it had to be measured for the army, find fault with Sicilian wheat, which it was lawful for him, if he wished, to measure out from the threshing-floor itself? So great a quantity of grain is given and exacted by your command. It is not enough; coins are commanded besides. They are given. It is too little. For the tithes of barley another sum of money is exacted. You order 30,000 sesterces to be given as gain. So from one city by force, by threats, by command, by the wrong-doing of the praetor, are wrung 33,000 medimni of wheat and besides 60,000 sesterces. But are these things obscure? Or, even if all men should wish, can they be obscure? Things which you did openly, which you commanded in the assize, which you compelled with all looking on; about which the Agyrine magistrates and chief five, whom you had summoned for the sake of your own gain, reported your acts and commands at home to their own senate. Their report has been entrusted to the public records by their laws. Their envoys, most noble men, are at Rome, who have said this very thing in testimony!

\ciceroSection{74} Hear the public letter of the Agyrines; then the public testimony of the city. Read. \textsc{Public letter, public testimony}. You have noted in this testimony, gentlemen, that Apollodorus, surnamed Pyragrus, the chief man of his city, weeping, bears witness and says that, since the name of the Roman people was heard and known by the Sicilians, never had the Agyrines spoken or done anything against any most lowly Roman citizen --- who now were being compelled to bear public testimony against a praetor of the Roman people, with great wrongs and great grief. By Hercules, Verres, your defence cannot stand against this one city; so great is the authority in those men's faithfulness, so great the grief in the wrong, so great the scruple in the testimony. Yet not one only, but whole cities, afflicted with like wrongs and hurts, pursue you with embassies and public testimonies.

\ciceroSection{75} Let us see now in turn how the city of Herbita --- honourable and once well-off --- was stripped and harassed by this man. And of what kind of men! The highest farmers, most far removed from the forum, the courts, the controversies, whom you ought to have spared and consulted for, most filthy man, and the kind of men of all most zealously to keep safe. In the first year the tithes of that field were sold for 18,000 modii of wheat. When Atidius (likewise this man's minister in tithe matters) had bought, and had come under the title of prefect to Herbita with the Venusians, and a place had been given him publicly to lodge in, the Herbitenses are compelled to give him as gain 38,800 modii of wheat, although the tithes had been sold for 18,000 modii. And this much gain are they compelled to give publicly at the time when already privately the farmers, despoiled and driven from their fields by the wrongs of the tithe-collectors, had fled away.

\ciceroSection{76} In the second year, when Apronius had bought the tithes for 25,800 modii of wheat and had himself come to Herbita with that troop and band of brigands, the people was compelled publicly to bring together for him as gain 21,000 modii of wheat and an addition of 4,000 sesterces. About the addition I am in doubt whether the wages were given to Apronius himself for his labour and shamelessness; but as to so great a number of wheat, who can doubt that it came to that grain-stealing brigand, like the Agyrine grain? But in the third year in this field he used a kingly custom. They say that the barbarian kings of the Persians and Syrians are wont to have several wives, and to assign cities to these wives in this manner: this city for the woman's headband, this for her necklace, this for her hair-ribbons. So they have whole peoples not only as accomplices of their lust, but even as ministers of it.

\ciceroSection{77} Learn that the licence and lust of this man, who said he was king of the Sicilians, was the same. Pipa is the wife of Aeschrio of Syracuse, whose name has been spread by this man's wickedness over all Sicily; about which woman very many verses were written above the tribunal and above the praetor's head. This Aeschrio, the imitation husband of Pipa, is set up as a new publican in the Herbitan tithes. The Herbitenses, when they saw that, if the price stuck at Aeschrio's bid, they would be despoiled at the will of a most lustful woman, bid up as far as they thought they could effect it. Aeschrio went above. For he was not afraid that under Verres as praetor a tithe-collecting woman could be afflicted with loss. The contract is knocked down at 8,100 medimni, almost half as much again as in the previous year. The farmers were utterly overthrown, and the more so because already in the previous years they had been afflicted and almost ruined. This man understood that the price had gone so high that no more could be wrung from the Herbitenses. He takes 600 medimni off the principal, and orders that in the books, for 8,100 medimni, 7,500 be entered.

\ciceroSection{78} The barley tithes of the same field Docimus had bought. This is the Docimus to whom this man had brought Tertia, daughter of Isidorus the mime, carried off by force from a Rhodian flute-player. Of this Tertia the authority with this man even more than that of Pipa, more than that of the rest, and almost as much as that of Chelidon at the city, prevailed in this man's Sicilian praetorship. There come to Herbita two rivals of the praetor, not troublesome --- the most dishonest representatives of the lowest little women. They begin to demand, to press, to threaten. They could not, however, although they wished, imitate Apronius. The Sicilians did not so dread Sicilians. When in every way, however, they were trumping up false charges, the Herbitenses promise bail to Syracuse. Once they have come there, they are compelled to give Aeschrio --- that is, Pipa --- as much as had been taken off the principal, 3,600 modii of wheat. He did not wish to give too much gain to the publican-woman from the tithes, lest perhaps from her nightly trade she should turn her mind to redeeming the revenues.

\ciceroSection{79} The Herbitenses thought the matter settled, when this man says, "What of the barley? Of Docimus, my dear little friend, what do you propose?" And he transacted this in his bedroom, gentlemen, and on his couch. They said no charge of that sort had been laid on them. "I do not hear: count out 41,000 sesterces." What were the wretches to do, or what to refuse? Especially when they saw on the couch the fresh tracks of the tithe-collecting woman, by which they understood him kindled to persistence. So this one allied and friendly city paid tribute under Verres as praetor to two of the foulest little women. And I now say that this number of grain and these sums of money were given to the tithe-collectors publicly by the Herbitenses; with which grain and with which sums they yet did not buy back their citizens from the wrongs of the tithe-collectors. For with the goods of the farmers already destroyed and plundered, this wage was being given to the tithe-collectors that some day they should depart from those men's fields and from the cities.

\ciceroSection{80} So, when Philinus of Herbita --- an eloquent and prudent man, noble at home --- spoke publicly about the calamity of the farmers, about their flight, about the fewness of those who remained, you noted, gentlemen, the groan of the Roman people, whose throng has never failed this case. About which fewness of the farmers I shall speak in another place. Now the thing I passed over does not seem to be wholly to be left aside. For, by the immortal gods, what he took off the principal --- in what manner indeed does it seem to you not only to be borne, but even to be heard?

\ciceroSection{81} One man only there has been since Rome was founded --- the immortal gods grant there be not another! --- to whom the commonwealth, compelled by the times and by domestic ills, gave herself wholly: Lucius Sulla. He could do so much that no man against his will could keep his goods, his fatherland, his life. He had so much spirit for audacity that he did not hesitate to say in a public meeting, when he was selling the goods of Roman citizens, that he was selling his own plunder. All his deeds we not only uphold but even, on account of fear of greater hurts and calamities, defend by public authority. This one thing has been reproached by several senatorial decrees: it was decreed that those to whom he had taken anything off the principal should refund the money to the treasury. The senate ruled that not even to him to whom it had granted all things was it lawful to lessen the totals of things made and acquired by the people.

\ciceroSection{82} If the senators judged that he could not remit anything from the sum due from the bravest men, will the senators judge that you rightly remitted to a most foul woman? He, about whom the Roman people had ordered by law that his own will could be to him in place of law, is yet in this one kind reproached by the religious obligation of the old laws: shall you, who were held entangled in all the laws, have wished that your lust should be to you in place of law? In him is reproached his having remitted out of money he had himself acquired: shall it be granted to you, who have remitted out of the principal of the revenues of the Roman people?

\ciceroSection{83} And in this kind of audacity even much more shamelessly did he act in the tithes of the people of Acesta. Having knocked them down to that same Docimus --- that is, to Tertia --- for 5,000 modii of wheat, and having added an addition of 1,500 sesterces, he compelled the Acestans to take publicly from Docimus that same amount; learn this from the public testimony of the Acestans. Read. \textsc{Public testimony}. You have heard for how much the city took the tithes from Docimus, for 5,000 modii of wheat and the addition. Learn now for how much he reported he had sold them. \textsc{The law of the tithes sold under Gaius Verres as praetor}. Under this head you see 3,000 modii of wheat taken off the principal, which, when he had taken from the sustenance of the Roman people, from the sinews of the revenues, from the blood of the treasury, he made over as a gift to Tertia the mime. Was it more shamelessly that he took from allies, or more foully that he gave to a harlot, or more dishonestly that he stole from the Roman people, or more audaciously that he altered the public records? From the strictness of these men shall any force snatch you, or any largesse? It shall not. But if it should snatch you, do you not understand that these things which I have long been speaking of belong to another inquiry --- to the trial for embezzlement?

\ciceroSection{84} So I shall reserve this whole kind of charge entire for myself: I shall return to the case I have set in hand, of the grain and the tithes. He, when he was laying waste the largest and most fertile fields by himself --- that is, by Apronius, the second Verres --- against the smaller cities had others whom he was wont to let loose like dogs, worthless and dishonest men, by whom he was compelling either grain or money to be given publicly. Aulus Valentius is in Sicily an interpreter, by whom this man as interpreter was wont to use not for the Greek tongue but for thefts and disgraces. This interpreter, a worthless and needy man, suddenly becomes a tithe-collector. He buys the tithes of the wretched and meagre field of Lipara for 600 medimni of wheat. The Liparenses are summoned. They themselves are compelled to take the tithes and to count out to Valentius 30,000 sesterces of gain. By the immortal gods, which will you take for your defence: that you sold the tithes for so much less that to the 600 medimni the city, of its own will, at once added 30,000 sesterces of gain (that is, 2,000 medimni of wheat); or, that, when you had sold the tithes for a great sum, you wrung this money out from the unwilling Liparenses?

\ciceroSection{85} But why do I ask of you what you mean to defend, rather than learn from the city itself what was done? Read out the public testimony of the Liparenses, then in what manner the coins were given to Valentius. \textsc{Public testimony. How the payment was made, from the public records}. Was even this so small a city, so far set apart from your hands and from your sight, sundered from Sicily, set on a slight and uncultivated island, when heaped with your other and greater wrongs, made to be your plunder and your gain in this kind of grain too? The whole island you had handed over to one of your fellow-revellers as if some little gift; from this also these grain-gains were exacted as if from inland places. So those who through so many years under you as praetor were wont to ransom their little fields from pirates --- these same men ransomed themselves at a price imposed by you.

\ciceroSection{86} What of the people of Tissa, a very small and slight city, but most painstaking and most thrifty farmers --- is there not snatched from them under the name of gain more than the whole grain they had ploughed up? Whom you sent Diognetus, a Venusian, as your tithe-collector, a new kind of publican. Why, on this man's authority, do not public slaves at Rome too go after the revenues? In the second year the Tissans are forced to give 21,000 sesterces of gain unwilling. In the third year they were compelled to give 12,000 modii of wheat as gain to Diognetus the Venusian. This Diognetus, who out of public revenues makes such gains, has no underling, no peculium at all. Doubt now if you can whether the Venusian attendant of this man took such a number of wheat for himself or exacted it for him.

\ciceroSection{87} Learn this from the public testimony of the Tissans. \textsc{Public testimony of the Tissans}. Obscurely, gentlemen, is the praetor himself the tithe-collector, when his attendants exact the grain from cities, command sums of money, and themselves carry off somewhat more gain than they will give to the Roman people under the name of tithes! This was the equity in your command, this the dignity of the praetor: that you wished Venus's slaves to be lords of the Sicilians. This was the choice, this the distinction under you as praetor: that the farmers were in the number of slaves, the slaves in the number of publicans.

\ciceroSection{88} What? The wretched Amestratines, when such great tithes had been imposed upon them that nothing was left for them, were they not yet compelled to count out money? The tithes are knocked down to Marcus Caesius, while the envoys of the Amestratines were present. At once Heraclius the envoy is compelled to count out 22,000 sesterces. What is this? What is this plunder, what is this force, what is this rape of allies? If a charge had been laid on Heraclius by the senate to buy, he would have bought; if there was none, how could he of his own accord count out money? He reports back to Caesius that he has given.

\ciceroSection{89} Learn the report from the public records. Read. \textsc{From the public records}. By what senatorial decree had this been entrusted to the envoy? By none. Why did he do it? He was compelled. Who says this? The whole city. Read out the public testimony. \textsc{Public testimony}. From the same city in the second year by like reasoning money was wrung out and given to Sextus Vennonius, as you have learned from the same testimony. But the Amestratines, slight men --- when you had sold their tithes to Bariobalus the Venusian for 800 medimni (learn the names of these publicans!) --- you compelled them to add more gain than they had been sold for, when they had been sold for a great sum. They give Bariobalus 850 medimni and 1,500 sesterces. Surely this man would never have been so mad as to suffer more grain to be paid to a Venusian slave than to the Roman people out of a field of the Roman people, unless all that plunder under the name of the slave came back to him himself.

\ciceroSection{90} The Petrines, when their tithes had been knocked down for a great sum, were yet most unwilling forced to give Publius Naevius Turpio --- a most dishonest man, who under Sacerdos as praetor was condemned for wrongs --- 52,000 sesterces. Did you sell the tithes so loosely that, when the medimnum was 15 sesterces, and the tithes had been sold for 3,000 medimni (that is, 45,000 sesterces), 52,000 sesterces of gain were given to the tithe-collector? But you sold the tithes of that field for a very great sum. Of course he is glorying in this: not that gain was given to Turpio, but that money was snatched from the Petrines.

\ciceroSection{91} What of the Halicyenses, whose dependents pay the tithes, while they themselves hold their own fields free of tribute --- were they not also compelled to give the same Turpio (when the tithes had been sold for 100 medimni) 15,000 sesterces? If you could prove what you most wish --- that these gains came to the tithe-collectors, that you yourself touched none of it --- yet these sums of money, taken and gathered by your force and wrong, ought to be your fault and the cause of your condemnation. But when you can persuade no one of this --- that you were so mad as to wish Apronius and Turpio, slave-men, to grow rich at the peril of yourself and your children --- do you suppose any will doubt that all this money was sought by those emissaries for you?

\ciceroSection{92} To Segesta also, a free city, the Venusian Symmachus is sent as tithe-collector. He brings letters from this man, that contrary to all senatorial decrees, contrary to all rights, contrary to the Rupilian law, the farmers should give him bail outside their forum. Hear the letters which he sent to the Segestans. \textsc{Letter of C. Verres}. How this Venusian made fools of the farmers, learn from one settlement of an honourable and well-favoured man; for the rest are of the same kind.

\ciceroSection{93} Diocles is a Panhormitan, surnamed Phimes, an illustrious and noble man. He was ploughing a field hired in Segestan territory; for there is no commerce in that field for any man. He had it on lease for six thousand sesterces. As tithe, when he had been beaten by the Venusian, he settled at 16,000 sesterces and 654 medimni. Learn this from his own books. \textsc{The entry of Diocles the Panhormitan}. To this same Symmachus, Gaius Annaeus Brocchus, a senator --- a man of that splendour, that virtue, of which all of you think --- was forced to give coins besides the grain. Did such a man, a senator of the Roman people, go to make gain for a Venusian slave under you as praetor?

\ciceroSection{94} If you did not reckon this order to surpass in dignity, did you not at least know this --- that it judges? Before, when the equestrian order judged, dishonest and rapacious magistrates in the provinces were the servants of the publicans; they distinguished those who were in their works; whatever Roman knight they had seen in the province, they followed up with kindnesses and liberality. Nor did that arrangement help the guilty so much as it harmed many, when they had done anything against the interest and will of that order. There was kept up at that time, somehow, by them, as if by common counsel, this principle diligently: that whoever had thought one Roman knight worthy of insult, should be judged worthy of trouble by the whole order:

\ciceroSection{95} you so despised the senatorial order, so levelled all things to your wrongs and your lusts, so held it set down and resolved in your mind, to challenge as judges all who lived in Sicily, or who had touched Sicily under you as praetor, that you did not even think of this: that you were going to come before judges of the same order? In whom, even if from their own private hurt no grief should sit, yet there would be that thought: that they had been despised in another's wrong, and the dignity of the order had been contemned and cast aside. Which by Hercules, gentlemen, seems to me not to be borne lightly. For an insult has a certain sting which modest men and good men can least bear.

\ciceroSection{96} You despoiled the Sicilians; for they are wont to be mute under their own wrongs. You harassed the businessmen; for they come to Rome unwilling and rarely. You handed over Roman knights to Apronius's wrongs; for what harm can they now do, who are not allowed to judge? What? When you afflict senators with the highest wrongs, what else are you saying than this: "Give me even that senator, that this most ample senatorial name may seem to be born not only for the envy of the unskilled, but even for the insult of the dishonest?"

\ciceroSection{97} Nor did he do this in the case of Annaeus alone, but in the case of all senators --- so that the name of the order weighed not so much for honour as for ignominy. Toward Gaius Cassius, a most distinguished and bravest man, when he was at that very time consul in this man's first year, he used such great wickedness that, when his wife --- a leading lady --- had her father's plough-lands at Leontini, he ordered the tithe-collectors to carry off all the grain. This man, Verres, you shall have as witness in this case, since you took care not to have him as judge.

\ciceroSection{98} You, however, gentlemen, ought to think there is something common and joined between us. Many duties have been laid on this order, many labours, many perils, not only of laws and trials, but even of rumours and of the times. So is this order set as it were exposed and lifted up on high, that it can seem to be blown round by all winds of envy. In so wretched and unjust a condition of life, shall we not even keep this, gentlemen: that we should not seem to our own magistrates, in upholding our own right, the most contemptible and most despicable of men?

\ciceroSection{99} The Thermitans sent men to buy the tithes of their own field. They thought it greatly in their public interest that they should be bought, however dear, rather than that they should fall in with one of this man's emissaries. There was set against them a certain Venuleius to buy. He did not cease bidding. They contended as long as it seemed it could in any way be borne. At last they ceased bidding. The contract is knocked down to Venuleius for 8,000 medimni of wheat. Posidorus the envoy reports back. Although this seemed to all not to be borne, yet Venuleius is given, lest he should come on, 7,000 modii of wheat and besides 2,000 sesterces. From which it easily appears what is the wage of the tithe-collector, what the praetor's plunder, seems to be. Give me the records of the Thermitans and the testimony. \textsc{Records of the Thermitans, and testimony}.

\ciceroSection{100} The Imacharenses, with all their grain now carried off, with all your wrongs now wrung out of them, you compelled, wretched and ruined as they were, to make a tribute, that they might give Apronius 20,000 sesterces. Read the decree on the tribute and the public testimony. \textsc{Senate-decree on the tribute being collected. Testimony of the Imacharenses}. The Hennenses, when the tithes of the Hennensian field had been sold for 8,200 medimni, were compelled to give Apronius 18,000 modii of wheat and 3,000 sesterces. Mark, please, what a quantity of grain is being exacted out of every tithe-paying field. For my speech runs through all the cities which owe tithes, and in this kind, gentlemen, I am now occupied --- not in that in which the farmers were one by one overthrown in all their goods, but in showing what gains were given publicly to the tithe-collectors, that they should some day depart full and gorged from those men's fields and cities, with this heap of their gain.

\ciceroSection{101} Why did you command the Calactines in the third year to give the tithes of their field, which they had been wont to give at Calacte, to Marcus Caesius the tithe-collector at Amestratus? --- which they had not done before you were praetor, nor had you yourself so laid down for two years before. Why was Theomnastus the Syracusan let loose by you onto the Mutycensian field? Who so harassed the farmers that they were compelled by want, for the second year's tithes (a thing I shall show in other cities also), to buy wheat of necessity.

\ciceroSection{102} Now from the agreements of the Hyblenses you will understand --- which were made with Gnaeus Sergius the tithe-collector --- that six times as much as had been sown was carried off from the farmers. Read the sowings and the agreements from the public records. Learn the agreements of the Menaeans with the Venusian slave. Learn likewise the declarations of sowings and the agreements of the Menaeans. Will you suffer it, gentlemen, that from allies, from the farmers of the Roman people, from those who labour for you, who serve you, who so wish the Roman commons to be fed by themselves that there may be left to themselves and their children only enough for them to be fed --- that from these by the highest wrong, by the bitterest insults, somewhat more was carried off than had been grown?

\ciceroSection{103} I feel, gentlemen, that I must now hold my speech in measure and shun your weariness. I shall not turn longer in one kind, and so shall I take away other things from my speech that I shall yet leave them in the case. You shall hear the complaints of the Agrigentines, bravest of men, most diligent of farmers. You shall learn the grief and wrongs of the Entellans, men of the highest labour and the highest industry. The hurts of the Heracleans, the Geloans, the Soluntines shall be set forth. You shall learn the fields of the Catinaeans --- most well-off and most friendly men --- harassed by Apronius. You shall understand that the Tyndaritan, a most noble city, the Cephaloeditan, the Haluntine, the Apolloniate, the Enguine, the Capitine were ruined by this iniquity of the tithes. To the Inenses, Murgentines, Assorines, Helorines, Ietines nothing at all was left. The Cetarini, the Scherini, men of small cities, were utterly cast down and ruined. In short, that all the tithe-paying fields through three years gave tribute to the Roman people in the tithe portion, and to Gaius Verres in all that remained --- and that to most farmers absolutely nothing was left over; but if to anyone anything was left or remitted, it was only as much as overflowed from the part with which his greed was satisfied.

\ciceroSection{104} I have left for myself the fields of two cities, gentlemen, almost the best and most noble: the Aetnensian and the Leontine. Of these fields I shall let go the gains of three years; I shall pick out one year, that I may the more easily set out what I have begun. I shall take the third year, both because it is the most recent and because it was so administered by this man that, when he saw he was certainly about to leave the province, he did not care if he was about to leave no farmer at all in Sicily. We shall deal with the tithes of the Aetnensian and Leontine fields. Mark, gentlemen, diligently. The fields are fertile; it is the third year; the tithe-collector is Apronius.

\ciceroSection{105} Of the Aetnensians I shall say very little; for they themselves spoke publicly in the previous hearing. You hold in memory that Artemidorus the Aetnensian, chief of that embassy, said publicly that Apronius came to Aetna with Venus's slaves; that he summoned the magistrates to him, ordered couches to be spread for him in the middle of the forum; that he was wont to feast daily not only in public but even out of the public funds; that, while at his banquets a band played and the largest cups were served, the farmers were wont to be detained, and from them, not only by wrong but even by insult, as much grain was wrung out as Apronius had commanded.

\ciceroSection{106} You have heard these things, gentlemen, all of which I now pass over and leave aside. I say nothing of Apronius's luxury, nothing of his insolence, nothing of the licence allowed him by this man, nothing of his singular wickedness and shamefulness: I shall speak only of the gain and profit of one field and one year, that you may the more easily make a guess of three years and of all Sicily. But the reckoning of the Aetnensians is brief for me; for they themselves came, they themselves carried down the public records. They taught you what little gain that no-bad man, the praetor's familiar Apronius, had made. Learn it, please, from their own testimony. Read. \textsc{Testimony of the Aetnensians}. What say you? Speak, speak, please, more clearly, that the Roman people may hear about its own revenues, its own farmers, its own allies and friends. 50,000 medimni, 50,000 sesterces, by the immortal gods! One field, in one year, gives Apronius 300,000 modii of wheat and besides 50,000 sesterces of gain! Were the tithes sold for so much less than they were worth, or, since they had been sold for sum great enough, was so great a quantity of grain and money carried off from the farmers by force? For whichever of these you say, in that the fault and the charge will stick.

\ciceroSection{107} For this you will not say --- and would that you did say it! --- that so much did not come to Apronius. So I shall hold you fast not only by the public agreements but by the private agreements and writings of the farmers, that you may understand that you have not been more diligent in committing thefts than I in catching them. Will you bear this? Will any defend it? Will these men, if they wish to rule otherwise about you, sustain it: that in one coming, out of one field, Quintus Apronius (besides that ready money I have spoken of) carried off 300,000 modii of wheat under the name of gain?

\ciceroSection{108} What? Do the Aetnensians say this alone? Indeed the Centuripines too, who hold by far the greatest part of the Aetnensian field, whose envoys --- the most noble men, Andron and Artemo --- the senate gave instructions in those things which publicly concerned their own city. Of those wrongs which the Centuripine citizens received not in their own but in another's borders, the senate and people of Centuripa would not send envoys. The Centuripine farmers themselves --- which is the greatest number in Sicily of most honourable and most well-off men --- chose three envoys, their own citizens, that by their testimony you might learn the calamities not of one field but of almost all Sicily. For the Centuripines plough almost throughout all Sicily; and they are weightier and surer witnesses against you, that the rest of the cities are moved by their own hurts alone, while the Centuripines, because they have possessions in the borders of almost all the cities, have felt the losses and hurts of the rest of the cities too.

\ciceroSection{109} But, as I said, the reckoning of the Aetnensians is sure, sealed in both public and private records. The pension of my diligence has rather to be exacted from the Leontine field, for this cause: that the Leontines themselves publicly have not much helped me. For under this man as praetor those wrongs of the tithe-collectors did not hurt them. Rather, gentlemen, they helped them. It will perhaps seem wonderful to you or incredible that, in such great hurts of the farmers, the Leontines, who were chiefs of the grain administration, were untouched by hurts and wrongs. The reason is this, gentlemen: that in the Leontine field, except for the one household of Mnasistratus, no Leontine possesses a clod. So you have heard the testimony of Mnasistratus, gentlemen, a most honourable and excellent man. Look not for the rest of the Leontines, whom not only Apronius could not hurt in their fields, but no storm even. For not only did they take no hurt, but in those Apronian rapines they were occupied with gain and profit.

\ciceroSection{110} Wherefore, since the Leontine city and embassy has failed me on the cause I have spoken of, I myself must take up the reckoning and find a way by which I may come to Apronius's gain --- nay, by which I may come to this man's huge and monstrous plunder. The tithes of the Leontine field, in the third year, were sold for 36,000 medimni of wheat, that is, 216,000 modii of wheat. For a great sum, gentlemen, a great sum. Nor can I deny it. So it is necessary that the tithe-collectors took either a loss or surely no great gain. For this is wont to be the case for those who have bought for a great sum.

\ciceroSection{111} What if I show that in this one purchase a hundred thousand modii of wheat were made as gain? What, two hundred thousand? What, three hundred thousand? What, four hundred thousand? Will you yet doubt for whom such a great plunder was sought? Some will say I am unjust, who from the size of the gain make a guess at the theft and plunder. What if I show, gentlemen, that those who make 400,000 modii of gain would have made a loss, had your iniquity, had your recoverers from your own staff, not intervened --- can anyone in such great gain and such great iniquity doubt that it was on account of your wickedness that you made such great gains, on account of the size of the gain that you wished to be wicked?

\ciceroSection{112} How then shall I attain to this, gentlemen, that I may know how much gain has been made? Not from Apronius's books, which when I sought them out I did not find, and when I had brought him into court I wrung from him a denial that he kept books. If he was lying, why was he taking them out of the way, if those books were not going to harm you? If he had compiled no records at all, does not that itself signify enough that he did not carry on his own business? For the manner of tithe-collectors is such that it cannot be carried on without very many writings; for the names of the farmers one by one, and the agreements with each, the tithe-collectors must necessarily pursue and write up. The acres all the farmers declared by your command and your institution. I do not think anyone declared less than he had ploughed, when so many crosses, so many punishments, so many recoverers from your staff were set out. In an acre of the Leontine field about a medimnum of wheat is sown, in continuous and even sowing. The field yields with the eighth, when it is well managed; but, with all the gods helping, with the tenth. If this happens at any time, then it comes about that there is just so much tithe as you have sown --- that is, that, however many acres have been sown, that many medimni of tithe are owed.

\ciceroSection{113} Things being so, first I say this: that the tithes of the Leontine field were sold for many thousands of medimni more than the thousands of acres sown in the Leontine field. But if it could not happen that they should plough up more than ten medimni from an acre, while a medimnum had to be given out of an acre as tithe, when the field --- which very rarely happens --- had brought forth with the tenth, what reckoning was there for the tithe-collectors (if indeed it was tithes, and not the goods of the farmers, that were being sold) to buy the tithes for somewhat more thousands of medimni than acres were sown? In the Leontine field the subscription and declaration of acres is no more than 30,000; the tithes were sold for 36,000 medimni. Did Apronius err, or rather was he mad? Nay, he would have been mad then, if it had been allowed the farmers to give what they owed, and not been necessary to give what Apronius had commanded.

\ciceroSection{114} If I show that no one gave less than three medimni an acre as tithe, you will grant, I suppose, that, since the field's fruit was reaped at the tenth, no one gave less than three tithes. And this was sought from Apronius as a kindness: that they might be allowed to settle at three medimni for each acre. For when from many four, even five, were being demanded, while from many not only no grain but not even chaff was being left out of all the harvest and out of a year's labour, then the Centuripine farmers (whose number in the Leontine field is the greatest) gathered into one place and chose Andron of Centuripa, a man chief among the honourable and noble men of his city, as envoy to Apronius (the same whom the Centuripine city has now sent at this time as envoy and witness to this trial), that he might plead the farmers' case before him, and seek of him not to exact more than three medimni per acre from the Centuripine farmers.

\ciceroSection{115} This was scarcely got from Apronius as a great kindness for those who were even then still uninjured. When this was being got, this in fact was being got: that, instead of single tithes, three tithes might be given. But if this had not been your business, they would have asked of you not to give more than single tithes, rather than of Apronius not to give more than three tithes. Now, that I may pass over for the moment those things which Apronius set down against the farmers in royal --- nay, tyrannical --- fashion, and that I may not call upon those from whom he snatched all the grain, and to whom he left nothing not only of the harvest but not even of their own goods --- of these three medimni, which he granted as a kindness and a favour, learn how much gain is made.

\ciceroSection{116} The declaration of the Leontine field is 30,000 acres; this comes to 90,000 medimni of wheat, that is, 540,000 modii. With 216,000 modii deducted (for which the tithes were sold), there remain 324,000 modii of wheat. Add three fiftieths of the whole sum of 540,000 modii; this comes to 32,400 modii of wheat (for from all there were besides being exacted three fiftieths). These are now towards 360,000 modii of wheat. But I had said 400,000 was made of gain. For I do not bring into this reckoning those for whom it was not allowed to settle for three medimni. Yet, that I may fill out the sum of my promise from this very reckoning, in addition to single medimni many were being compelled to give two sesterces, many one and a half sesterces of additional payment --- those who were paying the least, single sesterces apiece. To follow the least: since we have reckoned 90,000 medimni, let there be added by this new and worst precedent 90,000 sesterces.

\ciceroSection{117} Will he yet dare to tell me that he sold the tithes for a great sum, when from the same field he carried off himself half as much again as he sent to the Roman people? You sold the tithes of the Leontine field for 216,000 modii. If by law, for a great sum; if your lust was the law, if half-tithes were called tithes, you sold them for a small sum. For the yearly fruits of Sicily could have been sold for much more, had the senate or Roman people willed you to do that. For the tithes were often sold, when they were sold by the Hieronic law, for as much as now they were sold for under the Verrine law. Give me the tithes sold by Gaius Norbanus. \textsc{Tithes of the Leontine field sold by Gaius Norbanus}. And then no action was given on the measure of acres, nor was Artemidorus Cornelius a recoverer, nor was the Sicilian magistrate exacting from the farmer as much as the tithe-collector had declared, nor was a kindness sought from the tithe-collector, that they might settle at three medimni an acre, nor was the farmer compelled to give an addition of coins, nor to add three fiftieths in grain. And yet a great quantity of grain was being sent to the Roman people.

\ciceroSection{118} What do those fiftieths mean, what those additions of coins? By what right --- nay, by what custom --- did you do this? The farmer was giving coins. How, or whence? Who, if he had wished to be most generous, would have used a more heaped-up measure (as before they were wont to do in tithes, when they were sold on equal law and condition). He was giving coins! Whence? Out of his grain? As if under you as praetor he had had grain to sell. Out of the living, then, something had to be cut off, to be a source from which to Apronius, beyond those fruits of the plough-lands, this little crown of coins might be added. And further, did they give it willing or unwilling? Willing? They loved Apronius, I suppose. Unwilling? Why, unless they were compelled by force and trouble? Now this most senseless man in selling the tithes was making additions of coins to each tithe, and not little: he was adding two or three thousand sesterces. Through three years they amount to perhaps 500,000 sesterces. This he did neither by anyone's example nor by any right, nor did he refer the money to the treasury: nor will any man ever devise how he is to defend even this slight charge.

